% !TeX root = ../papers/nc-p3.tex
\section{第一長方形の水平全候補での文字枠とトーラス座標}
\label{sec:first-rectangle-horizontal-frames}
\begin{theorem}\label{thm:first-rectangle-horizontal-frames}
76型の水平台85で四次の次元が標準値であるための必要十分条件は $b+c\ne0$ である。
この全域において無印の指定始点の同型類は、次のトーラス開集合の各点で一意に表される。
\[
 (u,v,z)\in\mathbb G_m^3,\qquad
 (u+v)(1+uv)(z+v)(1-uz)(z-u)(1+vz)\ne0 .
\]
一位置移動は $J(u,v,z)=(v,1/u,-z)$。周期2の点は $u=v=1$ と $u=v=-1$ の二直線上に限られ、他の全点は周期4である。
全域で射影安定群は $(\mathbb Z/2)^3$、その中で全位置に可換な線型持上げを持つ四元部分群は一意である。
従ってこれらの同型・周期の記述は標準AS点を全て含む。さらに長さ5までの各経路空間は標準次元を持ち、
四次中心の文字別次元は $(3,2,0,0)$ で一定である。
\end{theorem}

\begin{proof}[四平面対を全域で復元する]
$b\ne1$ で $c=(t^2-b)/(1-b)$ と置く。定理\ref{thm:first-rectangle-orbits}の
四位置の対称性 $i\mapsto i\oplus e_p$, $e=(2,3,2,3)$ により、前節の部分traceと
固有空間の射影を挿入した縮約が与える分割は各位置で二対か一ブロックである。
目標より粗い候補は15個ある。少なくとも一位置に二対がある14候補に対し、各ブロックの隣接文字の
対角値の差を集め、台の非零因子および $b+c$ で飽和する。14個とも単位イデアルである。

全位置が一ブロックの場合には12個の部分traceの二次モーメントだけで足りる。
各位置で文字0と1の差の分子から単元を除くと、残りは全て $x=A$ の二次以下である。
係数を $1,x,x^2$ の順に置いた十二行三列行列の全三次小行列式を求める。
$\mathbb Q[b,t]$ で $bt(b-1)(t-b)(t+b)(t^2-b)$ を逆にすると、これらは単位イデアルを生成する。
同時消滅点があれば非零列 $(1,x,x^2)^t$ が核に入ることに矛盾する。
ここでは $b+c=(t^2-b^2)/(1-b)$ を使った。
決まった平面の四つ組 $(0,0,1,0),(1,1,0,1)$ に射影した反対切断には、いずれも
非零な二行二列に定数行列 $\left(\begin{smallmatrix}0&1\\1&0\end{smallmatrix}\right)$ がある。
前節と同じSegre交叉で四文字枠が復元される。
$b=1,c\ne1$ なら $t^2=1$ で、一位置移動により $b'=c\ne1$ に移る。
$b=c=1$ は定理\ref{thm:first-rectangle-orbits}の交線の論証で覆われる。
従って単項性に基づく指定始点の係数の一意性と安定群の結論は、ここでの全候補で成り立つ。
\path{computations/verify_first_rectangle_horizontal_frames.py} が14分割の飽和と十二行の行列の階数を検算する。
\end{proof}

\begin{proof}[基底曲面の一枚の座標]
$u=b/t,v=c/t$ とすると $t(1+uv)=u+v$ である。$b+c\ne0$ なら分母も非零である。
$z=((b+c)x-bc)/(bt)$ と置けば
\[
 t=\frac{u+v}{1+uv},\quad b=ut,\quad c=vt,\quad x=\frac{u(z+v)}{1+uv},\quad
 D=\frac{v(1-uz)}{u+v}.
\]
全非零係数の分子・分母の因子をまとめるとちょうど定理の開集合になる。
一位置移動の係数式に代入して $J$ を得る。$J^2(u,v,z)=(1/u,1/v,z)$ であり、
$J^2$ の固定条件と $u+v\ne0$ から $u=v=\pm1$。$z\ne0$ より奇数移動は固定点を持たない。
この二直線は $b=c=1,t=\pm1$ の特別成分との交線である。

二次関係の各色の像の基底は局所化上の単元小行列式で取れる。
$b\ne1$ の一枚および $b=1,t=\pm1$ の二枚で、それらを経路に順次挿入し、局所化の単元だけで掃き出す。
長さ0から5の各始点の階数は $1,4,10,20,35,56$ となる。同じ基底の四次・五次間の可換方程式の核の階数は
四文字順に $3,2,0,0$。追加の分母は要らず、これらの核と経路の余核は各ファイバーへの特殊化と可換である。
$b+c=0$ の排除は命題\ref{prop:parallel-antidiagonal-obstruction}による。
\end{proof}

\begin{corollary}
112型の水平主成分も $b+c\ne0$ では同じ $u,v$ を持ち、
$a=u(z+v)/(1+uv),\ y=u(w+v)/(1+uv)$ と書ける。左上行列式の非零条件は $zw\ne1$ である。
台147の一般無印商における指定始点の対合および一位置移動は
\[
 I(u,v,z,w)=(1/u,1/v,w,z),\qquad J(u,v,z,w)=(v,1/u,-z,-w).
\]
\end{corollary}
\begin{proof}
左上行列式を検出する $(b+c)ay-bc(a+y)+b^2(c-1)$ は
$u^2(u+v)^2(zw-1)/(1+uv)^3$ となる。残りは定理\ref{thm:parallel-horizontal-moduli}への代入である。
\path{computations/verify_horizontal_torus_coordinates.py} が両形の座標式と76型の全ファイバーに通用する低次の掃き出しを検算する。
\end{proof}
